% Part: sets-functions-relations
% Chapter: arithmetization
% Section: From N to Z

\documentclass[../../../include/open-logic-section]{subfiles}


\begin{document}

\olfileid{sfr}{arith}{int}
\olsection{Van $\Nat$ naar $\Int$: gehele getallen bouwen}
We beginnen met twee eenvoudige feiten:
	\begin{enumerate}
		\item Elk geheel getal kun je schrijven als $n - m$, met $n, m \in\Nat$.
		\item Dezelfde informatie uit $n - m$ kun je ook vastleggen met het geordende paar $\tuple{n, m}$.
	\end{enumerate}
We weten al dat de geordende paren van natuurlijke getallen de !!{element}s
van $\Nat^2$ zijn. We nemen ook aan dat we $\Nat$ al begrijpen. Samen
brengen die twee feiten ons op een eerste, na\"{i}ef idee: \emph{we behandelen gehele
getallen als geordende paren van natuurlijke getallen}.
Maar dat idee is te na\"{i}ef. We willen natuurlijk dat $0- 2 = 4 -
6$. Toch is $\tuple{0, 2 }\neq \tuple{4, 6}$: het zijn verschillende
paren. We kunnen dus niet zomaar zeggen dat $\Nat^2$ de verzameling
gehele getallen is.
Als we dit probleem algemener bekijken, zien we welke regel we nodig hebben:
	$$a - b = c - d \text{ precies dan als }a + d = c + b$$
(Dat dit de \emph{bedoeling} is, zie je door aan beide kanten $b$ en $d$ op
te tellen.) Daarom leggen we vast wanneer twee geordende paren bij
hetzelfde gehele getal horen. Zo krijgen we een equivalentierelatie. We noemen
die $\Intequiv$ en leggen haar zo vast:
$$\tuple{ a, b } \Intequiv \tuple{c, d}\text{ precies dan als }a + d = c + b$$
Nu moeten we laten zien dat dit echt een equivalentierelatie is.
\begin{prop} $\Intequiv$ is een equivalentierelatie.
\end{prop}
	\begin{proof}
	We moeten drie dingen over $\Intequiv$ laten zien: de relatie is
	reflexief, symmetrisch en transitief.
	\emph{Reflexiviteit: een paar hoort bij zichzelf.} Je ziet meteen dat
	$\tuple{a, b} \Intequiv \tuple{a, b}$, want $a + b = b + a$.
	\emph{Symmetrie: de relatie werkt ook andersom.} Stel dat $\tuple{a,
	b} \Intequiv \tuple{c, d}$. Dan is $a + d = c + b$. Dus is $c + b =
	a + d$, en daarmee $\tuple{c, d} \Intequiv \tuple{a, b}$.
	\emph{Transitiviteit: twee stappen kun je aan elkaar koppelen.} Stel
	dat $\tuple{a, b} \Intequiv \tuple{c, d}\Intequiv \tuple{m, n}$.
	Dan is $a + d = c + b$ en $c + n = m + d$. Daarom is $a + d + c +
	n = c + b + m + d$, en dus $a + n = m + b$. Hieruit volgt
	$\tuple{a, b } \Intequiv \tuple{m, n}$.
\end{proof}
Nu delen we de paren met deze relatie in groepen in. Zo'n groep heet
een equivalentieklasse:
\begin{defn}
		De gehele getallen zijn de equivalentieklassen van geordende paren
		van natuurlijke getallen onder $\Intequiv$. Dus $\Int =
		\equivclass{\Nat^2}{\Intequiv}$.
\end{defn}
Je kunt heel verschillend \emph{filosofisch} denken over deze
definitie, waarin we eenvoudig afspreken wat gehele getallen zijn.
Voordat we daarop ingaan, helpt het om eerst naar de technische kant te kijken.
Nu we weten wat de gehele getallen zijn, moeten we nog vastleggen hoe
optellen, vermenigvuldigen en vergelijken daarop werken. We schrijven
$\equivrep{m, n}{\Intequiv}$ voor de equivalentieklasse onder $\Intequiv$ waarin
$\tuple{m, n}$ een !!a{element} is.\footnote{Met de notatie uit
\olref[sfr][rel][eqv]{def:equivalenceclass} zouden we hiervoor
$\equivrep{\tuple{m,n}}{\Intequiv}$ schrijven. Dat leest alleen wat
lastiger.} Dit betekent:
	$$\equivrep{m, n}{\Intequiv} = \Setabs{\tuple{a, b} \in \Nat^2}{\tuple{a, b}\Intequiv \tuple{m, n}}$$
Nu leggen we optellen, vermenigvuldigen en vergelijken vast:
	\begin{align*}
		\equivrep{a, b}{\Intequiv} + \equivrep{c, d}{\Intequiv} &= \equivrep{a + c, b + d}{\Intequiv}\\
		\equivrep{a, b}{\Intequiv} \times \equivrep{c, d}{\Intequiv} &= \equivrep{a c + b  d, a  d + b c}{\Intequiv}\\
		\equivrep{a, b}{\Intequiv} \leq \equivrep{c, d}{\Intequiv} &\text{ precies dan als }a + d \leq b + c
	\end{align*}	
(Om de axioma's, de basisregels hierboven, gemakkelijker te lezen,
schrijf ik zoals gebruikelijk `$ab$' voor `$(a \times b)$'.) We moeten
nog controleren of deze definities echt doen wat ze \emph{moeten} doen.
Helemaal uitschrijven en narekenen kost nogal wat werk. Daarom staan de
details in \olref[check]{sec}. Kort gezegd: alles werkt.
Er blijft nog één punt over. We hebben de gehele getallen gebouwd met
natuurlijke getallen. Daardoor zijn de natuurlijke getallen \emph{zelf
geen gehele getallen}. In \olref[ref]{sec} komen we terug op wat dat
filosofisch betekent. Om ermee te rekenen hebben we in elk geval een
manier nodig om natuurlijke getallen \emph{als} gehele getallen te
gebruiken. Voor elke $n \in \Nat$ spreken we eenvoudig af dat $n_\Int
= \equivrep{n, 0}{\Intequiv}$. Nu moeten we nagaan of deze
afspraak goed werkt. Voor alle $m, n \in \Nat$ moet gelden:
\begin{align*}
	(m + n)_\Int &= m_\Int + n_\Int\\
	(m \times n)_\Int &= m_\Int \times n_\Int\\
	m \leq n &\liff m_\Int \leq n_\Int
\end{align*}
Die controles zijn niet moeilijk. Voor de tweede regel gebruiken we
gewoon de rekenregels voor natuurlijke getallen:
%\begin{proof}
%	Met de rekenregels voor natuurlijke getallen is meteen duidelijk:
	\begin{align*}
%		(m + n)_\Int  = \equivrep{m + n, 0}{\Intequiv} = \equivrep{m + n, 0 + 0}{\Intequiv} = \equivrep{m, 0}{\Intequiv} + \equivrep{n, 0}{\Intequiv} = m_\Int + n_\Int\\
		(m \times n)_\Int  &= \equivrep{m \times n, 0}{\Intequiv} \\
		&= \equivrep{m \times n + 0 \times 0, m \times 0 + 0 \times n}{\Intequiv} \\
		&= \equivrep{m, 0}{\Intequiv} \times \equivrep{n, 0}{\Intequiv} \\
		&= m_\Int \times n_\Int
	\end{align*}
%\end{proof}
De andere twee regels mag je zelf controleren.
\begin{prob}
	Laat zien dat $(m + n)_\Int = m_\Int + n_\Int$ en $m \leq n \liff
	m_\Int \leq n_\Int$ voor alle $m, n \in \Nat$.
\end{prob}
\end{document}
